\section{Consent, Divine Freedom, and the Analogical Judgment That God Loves}

When recoil loosens its grip, the creature does not return to the pre-reflective sleep of assuming it stands on its own. The arrest remains true: I am still wholly dependent right now on a Source that receives nothing back from me. What shifts is the posture of the will before that truth. \emph{Consent} is the rational creature's lucid, truthful ratification of its own received being---standing willingly inside the asymmetry without resentment, neither claiming aseity nor despising the existence it has been given.

Yet consent is not blind submission to inscrutable power. Looking steadily at the very asymmetry that first provoked vertigo---the fact that God needs nothing from any creature and gains nothing by creating---natural reason uncovers a decisive implication. Precisely because God lacks nothing, my present existence cannot be the product of divine compulsion, deficiency, or self-interest. To think that through to its root is to move from divine freedom, through divine justice and mercy, to the natural-theological judgment that God loves all existing things.

\subsection{Divine Freedom: Without Need, Without Compulsion}

Why does anything finite exist beside God? Two opposing errors have haunted philosophy whenever thinkers tried to answer that question.

The first error is \emph{necessitarianism} or emanationism: the claim that finite beings flow automatically out of the divine nature, the way heat radiates from fire or corollaries follow from a geometric theorem, so that God could no more refrain from creating than he could cease to be God. If that were true, then every creature would be a necessary byproduct of divine fullness rather than a willed gift---and worse, God would stand in necessary dependence on the world in order to realize his own goodness.

Aquinas refutes necessitarianism in \emph{ST} I, q.\ 19, a.\ 3: whether whatever God wills, he wills necessarily. He distinguishes \emph{absolute necessity} (which follows from the nature of the terms themselves, as human being necessarily implies animal) from \emph{necessity by supposition} (if Socrates is sitting, he necessarily sits while he is sitting, even though he was wholly free to stand). Applied to the divine will, this distinction settles the question. God's proper and principal object is his own infinite goodness, so God wills his own goodness by absolute necessity. Creatures, however, are willed only as ordered to God's goodness as their end---and an agent wills means to an end necessarily only when the end \emph{cannot be attained without them}, as a starving person must eat to live or a traveler must take a ship to cross the sea. Where an end already stands wholly complete without a given means, willing that means is free.

Aquinas states the conclusion directly:

\begin{quote}
Hence, since the goodness of God is perfect, and can exist without other things inasmuch as no perfection can accrue to Him from them, it follows that His willing things apart from Himself is not absolutely necessary. Yet it can be necessary by supposition, for supposing that He wills a thing, then He is unable not to will it, as His will cannot change.\endnote{Aquinas, \emph{ST} I, q.\ 19, a.\ 3, corpus. In Reply~4, he clarifies that this non-necessity reflects no deficiency or fickleness in God, ``but from a defect belonging to the nature of the thing willed, namely, that the perfect goodness of God can be without it; and such defect accompanies all created good.''}
\end{quote}

Patristic theology and solemn conciliar teaching affirm the same freedom. St.\ Augustine writes in \emph{City of God} XI.24 that ``God made what was made not from any necessity, nor for the sake of supplying any want, but solely from His own goodness, \emph{i.e.}, because it was good.''\endnote{Augustine, \emph{De Civitate Dei} XI.24 (Dods translation, 1871).} And the First Vatican Council defines in \emph{Dei Filius}, ch.\ 1, that God created ``not for the increase of His own happiness, nor to acquire but to manifest His perfection by the blessings which He bestows on creatures, with absolute freedom of counsel''---anathematizing in Canon~5 anyone who claims that God created by a necessity equal to that by which he loves himself.\endnote{First Vatican Council, \emph{Dei Filius}, ch.\ 1 and can.\ 5 \emph{De Deo rerum omnium creatore} (Manning 1912 English; Latin: \emph{non ad augendam suam beatitudinem, nec ad acquirendam, sed ad manifestandam perfectionem suam per bona, quae creaturis impertitur, liberrimo consilio}).}

\subsection{Free, Yet Not Arbitrary: Justice and Mercy at the Root of Being}

Rejecting necessitarianism raises the opposite fear: the second error, \emph{voluntarism} or arbitrary caprice. If God does not need creatures and does not create by necessity, is our existence merely the whim of an unaccountable sovereign who might act without reason or justice?

Aquinas closes that door in \emph{ST} I, q.\ 21, aa.\ 1 and~4. He first distinguishes two kinds of justice (\emph{ST} I, q.\ 21, a.\ 1). \emph{Commutative justice} governs mutual exchange and repayment between equals---buying, selling, paying wages. That kind of justice cannot belong to God, because no creature ever gives God something that was not already his; Aquinas cites St.\ Paul's blunt challenge from Romans~11:35: ``Who hath first given to Him, and recompense shall be made him?''\endnote{Aquinas, \emph{ST} I, q.\ 21, a.\ 1, corpus; cf.\ Rom.\ 11:35 (Douay-Rheims: ``Or who hath first given to him, and recompense shall be made him?'').}

Yet \emph{distributive justice}---the ordered wisdom by which a ruler gives to each member what befits its nature and station---belongs eminently to God. Because the object of will is good as apprehended by intellect, ``it is impossible for God to will anything but what His wisdom approves. This is, as it were, His law of justice, in accordance with which His will is right and just'' (\emph{ST} I, q.\ 21, a.\ 1, ad 2). God is never arbitrary, because his will is identical with his eternal wisdom. Once God freely chooses to create a human being, it is fitting to his wisdom that a rational soul should have hands, senses, and an ordered world suited to its flourishing---though even in paying each creature what befits its nature, ``He Himself is not the debtor, since He is not directed to other things, but rather other things to Him'' (\emph{ST} I, q.\ 21, a.\ 1, ad 3).

In \emph{ST} I, q.\ 21, a.\ 4, Aquinas traces divine justice to its deeper root in divine \emph{mercy}---understood broadly as the generous overcoming of defect, including the passage from non-existence into existence:

\begin{quote}
Now the work of divine justice always presupposes the work of mercy; and is founded thereupon. For nothing is due to creatures, except for something pre-existing in them, or foreknown. Again, if this is due to a creature, it must be due on account of something that precedes. And since we cannot go on to infinity, we must come to something that depends only on the goodness of the divine will---which is the ultimate end. We may say, for instance, that to possess hands is due to man on account of his rational soul; and his rational soul is due to him that he may be man; and his being man is on account of the divine goodness. So in every work of God, viewed at its primary source, there appears mercy.\endnote{Aquinas, \emph{ST} I, q.\ 21, a.\ 4, corpus. In Reply~4, answering the objection that creation presupposes nothing and therefore cannot involve justice or mercy, Aquinas states directly that justice is preserved in producing beings according to divine wisdom, ``And the idea of mercy, also, is preserved in the change of creatures from non-existence to existence.''}
\end{quote}

Every creaturely claim bottoms out in sheer gift. Nothing in non-existence had a right to be called into being; nothing in us earned the first instant of our existence, or earns the present instant either.

\subsection{The Analogical Judgment That God Loves All Existing Things}

If God needs nothing from creatures, creates with absolute freedom rather than compulsion, and holds every finite thing in existence right now by an act of wise will, the posture of that will toward the creature has a precise name. Aquinas sets out the argument in \emph{ST} I, q.\ 20, a.\ 2, resting its demonstration on natural metaphysics even as it opens under the \emph{Sed contra} of Wisdom~11:25 (``Thou lovest all things that are, and hatest none of the things which Thou hast made''):\endnote{Wis.\ 11:25--27 (Douay-Rheims): ``For thou lovest all things that are, and hatest none of the things which thou hast made: for thou didst not appoint, or make any thing hating it. And how could any thing endure, if thou wouldst not? or be preserved, if not called by thee? But thou sparest all: because they are thine, O Lord, who lovest souls.''}

\begin{quote}
God loves all existing things. For all existing things, in so far as they exist, are good, since the existence of a thing is itself a good; and likewise, whatever perfection it possesses. Now it has been shown above (Q.\ 19, A.\ 4) that God's will is the cause of all things. It must needs be, therefore, that a thing has existence, or any kind of good, only inasmuch as it is willed by God. To every existing thing, then, God wills some good. Hence, since to love anything is nothing else than to will good to that thing, it is manifest that God loves everything that exists. Yet not as we love. Because since our will is not the cause of the goodness of things, but is moved by it as by its object, our love, whereby we will good to anything, is not the cause of its goodness; but conversely its goodness, whether real or imaginary, calls forth our love, by which we will that it should preserve the good it has, and receive besides the good it has not, and to this end we direct our actions: whereas the love of God infuses and creates goodness.\endnote{Aquinas, \emph{ST} I, q.\ 20, a.\ 2, corpus. In Reply~4, Aquinas extends the same principle even to sinners: ``God loves sinners in so far as they are existing natures; for they have existence and have it from Him. In so far as they are sinners, they have not existence at all, but fall short of it; and this in them is not from God.''}
\end{quote}

This creative willing transforms the one-sided relation of \emph{ST} I, q.\ 13, a.\ 7. In human affairs, when another person does not need us, we tend to assume indifference, because finite love and need are routinely tangled together: we love a friend partly because their presence delights us, or a child partly because our hearts ache without them. Even at its purest, human love is \emph{responsive}---called forth by a goodness already standing before us that moves our appetite.

Divine love is the exact inverse. Because God has no lack, gains no perfection from creatures, and cannot be enriched by our existence, his willing our being cannot be an instrument of divine self-interest, insecurity, or appetite. Precisely because God needs nothing from me, the sole reason I exist at this second is that he freely wills my good---infusing and creating the very goodness by which I am. The asymmetry that first unsettled the intellect turns out to be the exact guarantee that our existence is pure gift rather than utility.

That judgment stands, however, under the strict discipline of \emph{analogy} (\emph{ST} I, q.\ 13, a.\ 5). Aquinas does not say that God loves ``just as we love,'' only louder or bigger. Univocal predication is impossible between God and creatures: in us, love is an accidental act distinct from our substance, often bound up with bodily emotion, attraction, and vulnerability, and every word we take from human life ``circumscribes and comprehends'' a finite perfection. Nor is ``love'' applied to God purely equivocally, as a mere homonym sharing nothing with our speech, for then nothing could be known or demonstrated about God from created things, against both philosophy and St.\ Paul's testimony in Romans~1:20. Instead, perfection-names are predicated of God and creatures \emph{analogically}---according to the creature's proportion to God as its cause, in whom all perfections pre-exist simply, unitedly, and in a mode surpassing every creaturely measure.\endnote{Aquinas, \emph{ST} I, q.\ 13, a.\ 5, corpus. Compare \emph{ST} I, q.\ 21, a.\ 1, ad 1: virtues bound up with sensitive passions cannot exist literally in God, ``since,\dots\ in God there are no passions,'' whereas acts belonging to the intellectual will---justice, liberality, love---belong properly and analogically to God.}

Thus the natural-theological judgment ``God loves all existing things'' is neither sentimental anthropomorphism nor an empty metaphor. It is a sober conclusion of reason: my present existence is neither an accident nor a debt paid to me nor a resource exploited by a needy deity, but the continuous effect of an uncreated, wise, entirely free will conferring good out of sheer plenitude.
